\chapter{Chapter 8 Refresher: Model Adaptation, LoRA, Quantization, and Deployment Reality --- Low-Rank Adaptation and Quantization}

\section*{Setting the Scene}
Deployment freezes most of \(\theta\) learned via Preface Step~8-style objectives and tunes cheap corrections under storage/power envelopes---approximation theory meets accounting.

\paragraph{Step A --- Low-rank hypothesis.}
\textbf{Problem.} Full matrices encode redundancy; updating all entries per task is wasteful.
\textbf{Insight.} SVD decomposes any matrix into a sum of rank-1 outer products ordered by singular value magnitude. Restrict \(\Delta W\) to rank-\(r\) manifolds---implicit claim task signals concentrate in few directions (Preface Step~6 thinks about singular values informally).
\textbf{Lineage.} Eckart--Young (1936) best low-rank approximation theorem; Hu et al.\ LoRA (2021) popularised modern adapter lore.
\textbf{Mental model.} Monitor singular spectrum drift when skeptics question rank budgets---a fast-decaying spectrum justifies small $r$; a flat spectrum does not.

\paragraph{Step B --- Quantisation as noisy channel.}
\textbf{Problem.} Continuous weights exceed SRAM budgets.
\textbf{Insight.} Rounding injects controlled distortion reminiscent of quantised communication channels---error scales with dynamic range misuse.
\textbf{Lineage.} Dettmers QLoRA et al.\ marry low-rank updates with int4 storage.
\textbf{Mental model.} Outliers behave like heavy-tail Preface Step~6 violations---scale grids adapt or suffer.

\paragraph{Step C --- Gate on measured drift.}
\textbf{Problem.} Approximations trade compute for risk.
\textbf{Insight.} Acceptance tests slice datasets like calibration bins---aggregate scores hide subgroup collapses tied to Preface Step~5 expectations.

Treat every shortcut as falsifiable: publish metrics or admit folklore.

\section*{Notation Introduced in This Chapter}
\begin{itemize}[leftmargin=1.5em]
  \item $W$: base weight matrix, $\Delta W$: adaptation update.
  \item $A,B$: low-rank factors in LoRA update.
  \item $d,k$: layer input/output dimensions.
  \item $r$: LoRA rank.
  \item $\alpha$: LoRA scaling constant.
  \item $\rho$: trainable parameter fraction versus full fine-tuning.
  \item $U,\Sigma,V$: left singular vectors, diagonal singular values, right singular vectors in SVD.
  \item $\sigma_i$: the $i$-th singular value (conventionally $\sigma_1 \ge \sigma_2 \ge \cdots \ge 0$).
  \item $\|M\|_F$: Frobenius norm, $\sqrt{\sum_{i,j}M_{ij}^2} = \sqrt{\sum_i \sigma_i^2}$.
\end{itemize}

\section*{What You Need To Recall Precisely}
\begin{itemize}[leftmargin=1.5em]
  \item \textbf{Singular Value Decomposition (SVD) and the low-rank approximation theorem.}

  Every real matrix $M \in \mathbb{R}^{d \times k}$ factors as
  \[
    M = U\Sigma V^\top,
  \]
  where $U \in \mathbb{R}^{d\times d}$ and $V \in \mathbb{R}^{k\times k}$ are orthogonal, and $\Sigma$ is diagonal with non-negative entries $\sigma_1 \ge \sigma_2 \ge \cdots \ge 0$. Expanding the product gives a sum of rank-1 outer products:
  \[
    M = \sum_{i=1}^{\min(d,k)} \sigma_i\, \mathbf{u}_i \mathbf{v}_i^\top.
  \]
  \emph{Eckart--Young theorem.} For any target rank $r$, the best rank-$r$ approximation in Frobenius norm (and in spectral norm) is obtained by keeping the top $r$ terms:
  \[
    M_r = \sum_{i=1}^{r} \sigma_i\, \mathbf{u}_i \mathbf{v}_i^\top,
    \qquad
    \|M - M_r\|_F = \sqrt{\sigma_{r+1}^2 + \cdots + \sigma_{\min(d,k)}^2}.
  \]
  No other rank-$r$ matrix is closer to $M$ in Frobenius norm. This is the theoretical foundation for using low-rank matrices as approximations.

  \emph{Connection to LoRA.} The LoRA hypothesis is that the task-relevant update $\Delta W$ has a rapidly decaying singular spectrum---most of its ``signal'' lives in a low-dimensional subspace. Restricting $\Delta W = \frac{\alpha}{r}BA$ (with $B \in \mathbb{R}^{d\times r}$, $A \in \mathbb{R}^{r\times k}$) is equivalent to parameterising the rank-$r$ subspace directly. Eckart--Young justifies this when the true update's singular values drop off quickly; when the spectrum is flat, a small $r$ will miss significant directions and underfitting follows.

  \emph{Diagnosing rank budget.} Inspect the singular value spectrum of $\Delta W$ after a full fine-tuning run (if affordable as a baseline). Rapid decay --- first few $\sigma_i$ capturing $\ge 90\%$ of $\|{\Delta W}\|_F$ --- supports a small LoRA rank. Slow decay flags either a hard task shift or a dataset with diverse sub-distributions, both arguing for larger $r$ or continued pretraining.

  \item \textbf{LoRA update form and intuition.}
  \[
    \Delta W = \frac{\alpha}{r}BA,\quad B\in\mathbb{R}^{d\times r},\ A\in\mathbb{R}^{r\times k}.
  \]
  Rank \(r\ll \min(d,k)\) forces \(\Delta W\) to live in an \(r\)-dimensional subspace of the full \(d\times k\) parameter space; optimisation restricts plasticity to directions where adaptation concentrates empirically, exchanging expressiveness for parameter savings.

  \item \textbf{Trainable parameter fraction:}
  \[
    \rho=\frac{r(d+k)}{dk}.
  \]
  \item \textbf{Quantization model:} map tensors through affine rescaling (per-channel or per-tensor scale/zero-point) into finite-bit grids---often asymmetric ranges capture outliers better than naive symmetric buckets. Quantisation noise concentrates where scales mismatch empirical tails (especially outliers), triggering subgroup drift unrelated to aggregate perplexity.
  \item \textbf{Error tradeoff:} memory/bandwidth gains versus approximation and calibration drift risk.
  \item \textbf{Budget accounting:} adaptation and base-model memory are separate components.
\end{itemize}

\section*{How These Results Build on Each Other}
\begin{enumerate}[leftmargin=1.5em]
  \item SVD decomposes any matrix into rank-1 components ordered by importance.
  \item Eckart--Young establishes that truncating to rank $r$ is the optimal approximation under Frobenius/spectral norm.
  \item The low-rank hypothesis applies this to weight updates: restrict $\Delta W$ to rank $r$ if task signal concentrates in few directions.
  \item LoRA parameterises the rank-$r$ subspace explicitly with factors $B, A$.
  \item Quantization adds a second approximation layer (bit-width reduction) with independent noise characteristics.
  \item Deployment gating measures drift from both approximations before release.
\end{enumerate}

\section*{Reminder Glossary (Term by Term)}
\begin{itemize}[leftmargin=1.5em]
  \item \textbf{Singular value decomposition (SVD):} factorisation $M = U\Sigma V^\top$ into orthogonal factors and a diagonal matrix of singular values; every real matrix has one.
  \item \textbf{Singular value ($\sigma_i$):} non-negative scalar measuring the ``size'' of the $i$-th rank-1 component; $\sigma_1 \ge \sigma_2 \ge \cdots$.
  \item \textbf{Singular spectrum:} the sequence $(\sigma_1, \sigma_2, \ldots)$; fast decay supports low-rank approximation, flat spectrum does not.
  \item \textbf{Eckart--Young theorem:} the best rank-$r$ approximation in Frobenius (or spectral) norm is the truncated SVD $M_r$; there is no better rank-$r$ matrix.
  \item \textbf{Frobenius norm:} $\|M\|_F = \sqrt{\sum_{i,j}M_{ij}^2} = \sqrt{\sum_i \sigma_i^2}$; natural measure of total approximation error.
  \item \textbf{Rank:} dimensionality of adaptation subspace; in LoRA, the number of singular directions parameterised explicitly.
  \item \textbf{LoRA scaling factor:} controls update magnitude relative to rank.
  \item \textbf{Quantizer:} mapping from continuous values to discrete levels.
  \item \textbf{Calibration set:} data used to estimate quantization scales.
  \item \textbf{Drift threshold:} maximum acceptable quality loss for deployment.
\end{itemize}

\section*{Worked Mini Example (SVD and Eckart--Young)}
Let $M = \begin{pmatrix}3&0\\0&1\\0&0\end{pmatrix}$. The SVD is $M = U\Sigma V^\top$ with
\[
  U = I_3,\quad \Sigma = \begin{pmatrix}3&0\\0&1\\0&0\end{pmatrix},\quad V = I_2,
\]
giving singular values $\sigma_1=3$, $\sigma_2=1$. The best rank-1 approximation keeps only $\sigma_1$:
\[
  M_1 = \begin{pmatrix}3\\0\\0\end{pmatrix}(1\ 0) = \begin{pmatrix}3&0\\0&0\\0&0\end{pmatrix},
  \qquad
  \|M - M_1\|_F = \sigma_2 = 1.
\]
No other rank-1 matrix achieves Frobenius error below 1. The ratio $\sigma_1^2/(\sigma_1^2+\sigma_2^2) = 9/10 = 90\%$ of total squared norm is captured by rank 1---a fast-decaying spectrum that supports the low-rank approximation.

\section*{Worked Mini Example (Parameter Fraction)}
Let $d=k=4096$ and $r=16$:
\[
\rho=\frac{16(4096+4096)}{4096\cdot4096}
=\frac{131072}{16777216}\approx 0.0078.
\]
This is about $0.78\%$ of full fine-tuning parameters for that layer.

\section*{Worked Example (Quality Gate Under Approximation)}
Suppose full fine-tuned baseline scores 78.4 on target benchmark.  
LoRA+quantized candidate scores 77.1 overall, but subgroup analysis shows:
\begin{itemize}[leftmargin=1.5em]
  \item common cases: -0.6
  \item rare long-context cases: -3.8
\end{itemize}
Reasoning: average drop may appear acceptable, but concentrated degradation violates deployment intent for long-context tasks.  
Decision: reject current quantization profile or introduce subgroup-aware gate.

\section*{Intuition Checkpoints}
\begin{itemize}[leftmargin=1.5em]
  \item SVD always exists and is unique (up to sign); it is not an approximation --- it is an exact factorisation. The approximation step is the truncation to rank $r$.
  \item Eckart--Young gives a lower bound on approximation error: $\|M-M_r\|_F \ge \sigma_{r+1}$. No representation trick can do better for a given rank.
  \item A low-rank update is a hypothesis about task-adaptation dimensionality, not a theorem. Inspect the singular spectrum to validate it.
  \item Quantization quality is distribution-dependent; outlier channels can dominate failure.
  \item ``Near parity'' must be defined metric-by-metric, not as a vague label.
\end{itemize}

\section*{Further Refresh}
\begin{itemize}[leftmargin=1.5em]
  \item Rank and linear-algebra background: see Chapter 4 refresher.
  \item Optimizer/memory arithmetic: see Chapter 2 refresher.
  \item Systems bottlenecks and throughput: see Chapter 7 refresher.
  \item \textbf{LoRA (Hu et al.)} --- low-rank adaptation mechanics.
  \item \textbf{QLoRA (Dettmers et al.)} --- quantization plus adapter training.
\end{itemize}
